
Il s'%
ensuit (V) que $x$ est la $m$-ème fonction symétrique élémentaire en les $\cl^{\bullet}(\underline{L}_r)-1$ ($r=1,\ldots,n$), donc est un polynôme homogène de degré $m$ à coefficients entiers dans les variables $(\cl^{\bullet}(\underline{L}_r)-1)$ ($r=1,\ldots,n$). Appliquant à nouveau l'hypothèse de récurrence et l'associativité de la loi (1.3.1) on est ainsi ramené à démontrer que, pour tout $\underline{O}_X$-module inversible $\underline{L}$, on a
\begin{equation}\tag{1.3.6}
(\cl^{\bullet}(\underline{L})-1)\cl_\bullet(\underline{O}_X)\in\Filt_{j-1}(X)
\end{equation}
Le schéma $X$ étant intègre, $\underline{L}$ et $\underline{O}_X$ sont isomorphes à des sous-modules du $\underline{O}_X$-module constant $\underline{R}_X$ des fonctions rationnelles sur $X$. Le $\underline{O}_X$-module somme $\underline{F}=\underline{O}_X+\underline{L}$ est cohérent, et on trouve
\begin{align*}
(\cl^{\bullet}(\underline{L})-1)\cl_\bullet(\underline{O}_X)
&=\cl_\bullet(\underline{L})-\cl_\bullet(\underline{O}_X)\\
&=(\cl_\bullet(\underline{F})-\cl_\bullet(\underline{O}_X))-(\cl_\bullet(\underline{F})-\cl_\bullet(\underline{L}))\\
&=\cl_\bullet(\underline{F}/\underline{O}_X)-\cl_\bullet(\underline{F}/\underline{L}).
\end{align*}
Mais les fibres au point générique des faisceaux $\underline{F}/\underline{O}_X$ et $\underline{F}/\underline{L}$ sont nulles, donc leurs supports dans $X$ sont rares ($X$ étant irréductible), donc de dimension $\leq j-1$. L'assertion (1.3.6) en découle immédiatement.

II. Reste à démontrer le cas général. On introduit le fibré en drapeaux
\[
g:D\longrightarrow X
\]
de $\underline{E}$ sur $X$, où on peut trouver des $\underline{O}_D$-modules inversibles $\underline{L}_1,\ldots,\underline{L}_n$ tels que
\begin{equation}\tag{1.3.7}
g^*(\cl^{\bullet}(\underline{E}))=\cl^{\bullet}(L_1)+\cdots+\cl^{\bullet}(\underline{L}_n)\in K^{\bullet}(D).
\end{equation}
Soient $e$ le point générique de $X$, $\underline{k}(e)$ le corps résiduel de $e$. Alors on peut trouver un point rationnel $u$ de la fibre générique $D_e$ de $D$ 
